\documentclass[12pt]{article}
\usepackage{amsfonts}
\usepackage{amsmath, amssymb, amsthm}
\textheight=21.4 cm \textwidth=15cm  \topmargin=-0.2 in
\oddsidemargin 0cm \evensidemargin -0.2cm
\parindent 24pt

\newtheorem{theorem}{Theorem}[section]
\newtheorem{lemma}[theorem]{Lemma}
\newtheorem{proposition}[theorem]{Proposition}
\newtheorem{corollary}[theorem]{Corollary}
\theoremstyle{definition}
\newtheorem{definition}[theorem]{Definition}
\theoremstyle{remark}
\newtheorem{remark}[theorem]{Remark}
\newtheorem{question}[theorem]{Question}

\newcommand{\cb}{\mathrm{cb}}
\newcommand{\CB}{\mathcal{CB}}
\newcommand{\B}{\mathcal{B}}
\newcommand{\K}{\mathcal{K}}
\newcommand{\T}{\mathcal{T}}
\newcommand{\Hil}{\mathcal{H}}
\newcommand{\Mn}{M_n}
\newcommand{\otmin}{\otimes_{\min}}
\newcommand{\ri}{\rightarrow}

\title{Injectivity, exactness and the weak expectation property\\ for higher duals of operator spaces}
\author{Yafei Zhao}
\date{}

\begin{document}
\maketitle

\begin{center}
\begin{minipage}{11.5cm}
{\bf Abstract:} We study how exactness, local reflexivity, the weak
expectation property (WEP), injectivity and nuclearity transfer between an
operator space $V$ and its second dual $V^{**}$. Known injectivity and
nuclearity laws are due to Dong and Tao and are cited rather than reproved.
We record elementary descent, show that exactness, local reflexivity and the
WEP do not admit simple bidual lift laws, and isolate the reliable substitute
equivalences given by weak$^*$ exactness, bidual injectivity, and the local
lifting property. The compact operators $\K(\Hil)$ and the trace class
$\T(\Hil)$ provide separating examples.
\par
{\bf Keywords:} operator space, second dual, exactness, injectivity,
nuclearity, weak expectation property, local reflexivity
\par
{\bf 2020 MR Subject Classification:} Primary 46L07; Secondary 47L25
\end{minipage}
\end{center}

\section{Introduction}

The local theory of operator spaces quantises the local theory of Banach
spaces \cite{Ru89,Ble92}. For an operator space $V$ one has the implications
\begin{equation}\label{eq:chain}
V\ \text{nuclear}\ \Longrightarrow\ V\ \text{exact}\ \Longrightarrow\ V\
\text{locally reflexive},
\end{equation}
the first due to Pisier \cite{Pi92} and the second to Effros, Ozawa and Ruan
\cite{EOR}. The same authors characterised nuclearity through the bidual:
$V$ is nuclear if and only if $V$ is locally reflexive and $V^{**}$ is
injective \cite{EOR}. For $C^*$-algebras this specialises, via Connes
\cite{Co76} and Choi--Effros \cite{CE76,CE77}, to ``$\mathcal A$ nuclear
$\iff\mathcal A^{**}$ injective''. Dong and Ruan \cite{DR} completed the
picture for exactness: $V$ is exact if and only if $V$ is locally reflexive
and $V^{**}$ is weak$^*$ exact.

Each of these results reads a global property of $V$ off the bidual. Here we
ask the reverse question: for a local property $P$, when does $V^{**}$ have
$P$, expressed through $V$ and a transfer condition? Dong and Tao \cite{DT}
settled injectivity and nuclearity; we cite their identities when needed and
focus on exactness, local reflexivity and the WEP. We first record elementary
descent (Section~\ref{sec:descent}), then show that for the latter three
properties no simple bidual law of the form ``$V^{**}$ has $P\iff V$ has $P$
plus a single transfer condition'' holds (Section~\ref{sec:nolift}). The trace
class $\T(\Hil)$ and the compact operators $\K(\Hil)$ provide the separating
examples.

\section{Preliminaries}

We follow \cite{ER00,Pi03} for operator space theory. For operator spaces
$V,W$ and a linear map $\varphi:V\ri W$, the amplification
$\varphi_n=\mathrm{id}_{\Mn}\otimes\varphi:\Mn(V)\ri\Mn(W)$ acts entrywise,
and $\|\varphi\|_\cb=\sup_n\|\varphi_n\|$. We write $\CB(V,W)$ for the
operator space of completely bounded maps, $V^*=\CB(V,\mathbb C)$ for the
dual, and $\iota_V:V\ri V^{**}$ for the canonical map; $\iota_V$ is a complete
isometry \cite[\S 3.2]{ER00}. We use the injective and projective operator
space tensor products $\otmin$ and $\widehat\otimes$ \cite{Ble92}.

Throughout, the five \emph{local properties} of an operator space $V$ are
those of \eqref{eq:chain} together with injectivity and the WEP.

\begin{definition}\label{def:props}
Let $V$ be an operator space.
\begin{itemize}
\item[(i)] $V$ is \emph{injective} if for every inclusion $X\subseteq Y$ of
operator spaces, each complete contraction $X\ri V$ extends to a complete
contraction $Y\ri V$.
\item[(ii)] $V$ is \emph{locally reflexive} if for every finite-dimensional
operator space $L$, each complete contraction $L\ri V^{**}$ is the
point-weak$^*$ limit of a net of complete contractions $L\ri V$;
equivalently $\CB(L,V^{**})=\CB(L,V)^{**}$ for every such $L$
\cite[Thm.~14.3.1]{ER00}.
\item[(iii)] $V$ is \emph{exact} if for every finite-dimensional $L\subseteq
V$ and $\varepsilon>0$ there are $n$ and $S\subseteq\Mn$ with
$d_\cb(L,S)<1+\varepsilon$ \cite{Pi92,Pi03}.
\item[(iv)] $V$ is \emph{nuclear} if the identity of $V$ factors approximately
through matrix algebras by complete contractions in the point-norm topology
\cite{EOR}.
\item[(v)] $V$ has the \emph{weak expectation property} (WEP) if for one
(equivalently every) completely isometric inclusion $\tau:V\hookrightarrow
\B(\Hil)$ there is a complete contraction $P:\B(\Hil)\ri V^{**}$ with
$P\circ\tau=\iota_V$ \cite[\S 15]{ER00}.
\end{itemize}
\end{definition}

We record the inclusions we use. Exactness passes to subspaces
\cite[\S 17]{Pi03}; injectivity passes to completely contractively
complemented subspaces \cite[\S 6.1]{ER00}; and an operator space which is
injective and dual as a Banach space has the form $eR(1-e)$ for an injective
von Neumann algebra $R$ and a projection $e\in R$ \cite{EOR}. We also use that
every dual operator space $V^*$ is completely contractively complemented in
its bidual $V^{***}$ by the adjoint of $\iota_V$, since
$(\iota_V)^*:V^{***}\ri V^*$ satisfies $(\iota_V)^*\circ\iota_{V^*}
=\mathrm{id}_{V^*}$ \cite[\S 3.2]{ER00}.

\section{Descent of local properties}\label{sec:descent}

We first observe that all the local properties in question pass from the
second dual down to the space, by the completely isometric inclusion
$\iota_V:V\hookrightarrow V^{**}$.

\begin{proposition}\label{prop:descent}
Let $V$ be an operator space.
\begin{itemize}
\item[(i)] If $V^{**}$ is exact, then $V$ is exact.
\item[(ii)] If $V^{**}$ is injective, then $V$ has the WEP.
\end{itemize}
\end{proposition}

\begin{proof}
(i) Exactness passes to subspaces, and $V$ embeds completely isometrically in
$V^{**}$ via $\iota_V$.

(ii) Fix a completely isometric inclusion $\tau:V\hookrightarrow\B(\Hil)$. The
complete contraction $\iota_V:V\ri V^{**}$ then extends, by injectivity of
$V^{**}$ applied to $\tau$, to a complete contraction $P:\B(\Hil)\ri V^{**}$
with $P\circ\tau=\iota_V$. Hence $V$ has the WEP.
\end{proof}

For injectivity and nuclearity, the corresponding lift laws and their
converses are due to Dong and Tao \cite{DT}; we use them only by citation.
For example, $V^*=\B(\ell^2)$ is injective for $V=\T(\ell^2)$, yet
$V^{***}$ is not injective \cite[after Prop.~2.7]{DT}, so injectivity does
not lift without an exactness hypothesis.

\section{Alternative equivalences and the failure of simple lifting}
\label{sec:nolift}

Dong and Tao \cite{DT} show that injectivity and nuclearity lift to higher
duals with transfer condition ``exact''. One might hope for similar bidual
laws for exactness, local reflexivity and the WEP. Proposition~\ref{prop:descent}
gives only descent. Reverse lifting fails, and the search for equivalences
leads to the following three substitutes.

\subsection{The Weak Expectation Property (WEP)}

For the WEP, the bidual property collapses to injectivity on the dual space
$V^{**}$ via a standard argument.

\begin{proposition}\label{prop:wep-equiv}
For any operator space $V$, the following are equivalent:
\begin{enumerate}
\item[(i)] $V^{**}$ has the WEP;
\item[(ii)] $V^{**}$ is injective.
\end{enumerate}
Consequently, if $V$ is locally reflexive, then
\[
        V^{**}\in\mathrm{WEP}\quad\Longleftrightarrow\quad V\ \text{is nuclear}.
\]
In particular, the same equivalence holds for every exact operator space $V$.
\end{proposition}

\begin{proof}
It is a well-known fact that for dual operator spaces, the WEP coincides with
injectivity. Indeed, (ii) $\Rightarrow$ (i) is trivial. Conversely,
let $X=V^{**}$, which is a dual operator space. If $X$ has the WEP, then for
any completely isometric inclusion $X\subset \B(\Hil)$, the canonical map
$\iota_X: X \ri X^{**}$ extends to a complete contraction $P:\B(\Hil)\ri X^{**}$.
Since $X$ is a dual space, the canonical projection $\pi:X^{**}\ri X$ is
completely contractive. The composition $\pi\circ P:\B(\Hil)\ri X$ is then a
completely contractive projection onto $X$, so $X$ is injective. 
The final assertion follows from the Effros--Ozawa--Ruan theorem:
$V$ is nuclear if and only if $V$ is locally reflexive and $V^{**}$ is
injective \cite[Thm.~4.5]{EOR}.
\end{proof}

Thus a WEP law for $V^{**}$ is not governed by the WEP of $V$ alone. It is
governed by injectivity of $V^{**}$, and by nuclearity of $V$ once local
reflexivity is assumed. The local reflexivity hypothesis cannot be omitted:
Kirchberg's example, quoted in \cite{EOR}, gives a separable non-nuclear
operator space $V$ with $V^{**}=\prod_{n=1}^{\infty}M_n$, hence with
$V^{**}$ injective and WEP.

\subsection{Exactness and Weak$^*$ Exactness}

Exactness does not lift to the second dual, even for highly tractable spaces.

\begin{theorem}\label{thm:exact-nolift}
There is an exact operator space $V$ with $V^{**}$ not exact.
\end{theorem}

\begin{proof}
Take $V=\K(\ell^2)$, the compact operators. Then $V$ is nuclear, hence exact
by \eqref{eq:chain}. Its bidual $V^{**}=\B(\ell^2)$ is not exact
\cite[\S 9.3]{Pi03}. 
\end{proof}

For $C^*$-algebras, Wassermann proved that a von Neumann algebra is exact if and
only if it is subhomogeneous \cite{Wa94}. Thus, $V^{**}$ being exact forces extremely
restrictive finite-dimensional-like constraints on $V$ in the self-adjoint case.
For arbitrary operator spaces the honest bidual equivalence is obtained by
applying Dong--Ruan's theorem to $V^{**}$:
\[
V^{**}\ \text{exact}\quad\Longleftrightarrow\quad
V^{**}\ \text{locally reflexive and}\ V^{****}\ \text{weak}^*\text{ exact}.
\]
This is a genuine equivalence, but it is not a condition on $V$ alone.
Dong--Ruan's original theorem also gives the balanced substitute
\begin{equation}
V\ \text{exact}\quad\Longleftrightarrow\quad V\ \text{locally reflexive and}\
V^{**}\ \text{weak}^*\text{ exact}.
\end{equation}

\subsection{Local Reflexivity}

Local reflexivity similarly fails to lift.

\begin{theorem}\label{thm:lr-nolift}
There is a locally reflexive operator space $V$ with $V^{**}$ not locally reflexive.
\end{theorem}

\begin{proof}
Again take $V=\K(\ell^2)$. As a nuclear operator space it is locally reflexive
by \eqref{eq:chain}. Its bidual $\B(\ell^2)$ is injective, but not nuclear.
Since an operator space $X$ is nuclear if and only if $X$ is locally reflexive and
$X^{**}$ is injective \cite{EOR}, applying this to $X=\B(\ell^2)$ shows that $\B(\ell^2)$
cannot be locally reflexive.
\end{proof}

The available equivalence is therefore the defining tensor-bidual criterion
applied to $V^{**}$:
\[
V^{**}\ \text{locally reflexive}
\quad\Longleftrightarrow\quad
\CB(L,V^{****})=\CB(L,V^{**})^{**}
\]
completely isometrically for every finite-dimensional operator space $L$
\cite[Thm.~14.3.1]{ER00}. This is precise but too tautological to be a
publishable transfer law. The example $V=\K(\ell^2)$ shows that adding only
local reflexivity of $V$ cannot suffice.

\subsection{A positive neighbouring result: LLP}

The local lifting property provides a useful contrast. Dong proved that
\[
        V^{**}\ \text{has LLP}
        \quad\Longleftrightarrow\quad
        V\ \text{has LLP and}\ V^*\ \text{is exact}
\]
for every operator space $V$ \cite{DongLLP}. This is the kind of bidual
equivalence one might hope for in the present problem, but it belongs to LLP,
not to exactness or local reflexivity. It also explains why replacing WEP by
LLP leads to a substantially stronger statement.

\section{Summary}

Injectivity and nuclearity lift to higher duals with transfer condition
``exact'' \cite{DT}. Exactness, local reflexivity and the WEP descend from
$V^{**}$ to $V$ but do not lift in the same form. The exactness line has a
true higher-bidual reformulation through weak$^*$ exactness \cite{DR}; the
local reflexivity line has only the tensor-bidual criterion of
\cite[Thm.~14.3.1]{ER00}; and the WEP line reduces to injectivity of
$V^{**}$, hence to nuclearity of $V$ under local reflexivity
(Proposition~\ref{prop:wep-equiv}). The positive theorem in this neighbourhood is
Dong's LLP result, where
$V^{**}$ has LLP if and only if $V$ has LLP and $V^*$ is exact \cite{DongLLP}.

\begin{thebibliography}{99}

\bibitem{Ble92} D.~P. Blecher, Tensor products of operator spaces II,
\emph{Canad. J. Math.} \textbf{44} (1992), 75--90.

\bibitem{CE76} M.-D. Choi and E.~G. Effros, Separable nuclear $C^*$-algebras
and injectivity, \emph{Duke Math. J.} \textbf{43} (1976), 309--322.

\bibitem{CE77} M.-D. Choi and E.~G. Effros, Nuclear $C^*$-algebras and
injectivity: the general case, \emph{Indiana Univ. Math. J.} \textbf{26}
(1977), 443--446.

\bibitem{Co76} A.~Connes, Classification of injective factors,
\emph{Ann. of Math.} \textbf{104} (1976), 73--115.

\bibitem{DR} Z.~Dong and Z.-J. Ruan, Weak$^*$ exactness for dual operator
spaces, \emph{J. Funct. Anal.} \textbf{253} (2007), 373--397.

\bibitem{DongLLP} Z.~Dong, On the local lifting properties of operator spaces,
\emph{Canad. J. Math.} \textbf{61} (2009), 1262--1278.

\bibitem{DT} Z.~Dong and J.~Tao, On the nuclearity for the dual operator
space, \emph{Proc. Indian Acad. Sci. (Math. Sci.)} \textbf{120} (2010),
105--111.

\bibitem{EOR} E.~G. Effros, N.~Ozawa and Z.-J. Ruan, On injectivity and
nuclearity for operator spaces, \emph{Duke Math. J.} \textbf{110} (2001),
489--521.

\bibitem{ER00} E.~G. Effros and Z.-J. Ruan, \emph{Operator Spaces}, London
Math. Soc. Monographs, New Series \textbf{23}, Oxford Univ. Press, 2000.

\bibitem{Pi92} G.~Pisier, Exact operator spaces, in \emph{Recent Advances in
Operator Algebras}, Astérisque \textbf{232} (1995), 159--186.

\bibitem{Pi03} G.~Pisier, \emph{Introduction to Operator Space Theory}, London
Math. Soc. Lecture Note Series \textbf{294}, Cambridge Univ. Press, 2003.

\bibitem{Ru89} Z.-J. Ruan, Subspaces of $C^*$-algebras,
\emph{J. Funct. Anal.} \textbf{76} (1988), 217--230.

\bibitem{Wa94} S. Wassermann, \emph{Exact $C^*$-algebras and related topics},
Seoul National University, 1994.

\end{thebibliography}

\end{document}
